\begin{ouverture}[college3]{13}{Circuits : schématiser et protéger}{Partie 3 · L’énergie, ses transferts et ses conversions}{Deux lampes peuvent-elles fonctionner indépendamment ?}
\objectif{Réaliser un schéma avec les symboles usuels.}
\objectif{Distinguer série, dérivation, branche et nœud.}
\objectif{Expliquer le risque d’un court-circuit et le rôle d’une protection.}
\prerequis{Boucle fermée, conducteur et isolant.}
\end{ouverture}
\pcPage{Le cours}{Circuits : schématiser et protéger}
\phantomsection\label{ess-13}
\essSection{1}{Un modèle du circuit}
Un générateur alimente un circuit ; des récepteurs convertissent l’énergie reçue. Pour un courant continu conventionnel, on représente le sens du courant de la borne positive vers la borne négative du générateur dans le circuit extérieur. Un interrupteur ouvert interrompt une branche. Les fils sont tracés avec une règle et les symboles ne représentent pas l’apparence photographique des objets.\par
\pcFigure{figures/c13-circuits.pdf}{À gauche : deux lampes en série. À droite : deux branches en dérivation sur le même générateur.}
\essSection{2}{Série et dérivation}
Deux dipôles sont en série lorsque le même courant les traverse dans une branche sans séparation entre eux. En dérivation, ils sont reliés aux mêmes deux nœuds et appartiennent à des branches différentes. Une branche est une portion entre deux nœuds ; un nœud relie au moins trois fils dans le modèle scolaire. Croiser des traits sans point de connexion ne signifie pas toujours les relier : la convention du schéma doit être lisible.\par
\clearpage
\essSection{3}{Le court-circuit}
Un fil de très faible résistance reliant directement les bornes d’un dipôle le court-circuite. Court-circuiter un générateur peut entraîner un courant très élevé, un échauffement et une détérioration. Même une pile peut chauffer dangereusement. On étudie le défaut sur schéma ou avec un matériel protégé sous conduite du professeur, jamais en reliant volontairement les bornes d’une pile par un fil.\par
\pcFigure{figures/c13-protection.pdf}{Schéma de principe : une protection ouvre le circuit en cas de surintensité. Ce dessin n’est pas un plan de câblage du secteur.}
\essSection{4}{Des protections complémentaires}
Un fusible ou un disjoncteur approprié peut interrompre une surintensité. L’isolation, la mise à la terre et la protection différentielle répondent à d’autres fonctions complémentaires dans une installation réelle. Pour les expériences du livre, utiliser exclusivement le matériel scolaire de très basse tension autorisé. Aucun montage n’est réalisé sur une prise domestique, même à partir d’un schéma simplifié.\par
\begin{pcEncadre}{method}{Un exemple expliqué}Deux lampes sont branchées sur deux branches distinctes entre les bornes du générateur. Ouvrir l’interrupteur placé uniquement dans la branche de la première coupe ce trajet ; l’autre peut rester alimentée si son circuit est intact. En série, un interrupteur qui ouvre l’unique boucle coupe le courant dans les deux lampes.\end{pcEncadre}
\pcPage{Expérimenter et retenir}{Circuits : schématiser et protéger}
\phantomsection\label{activite-13}
\begin{pcEncadre}{method}{Investigation en classe}\textbf{Cahier des charges.} Réaliser avec le matériel scolaire deux lampes dont l’une peut être commandée sans interrompre l’autre. Dessiner le schéma avant de brancher. Faire vérifier le choix de la tension et les connexions par le professeur. Tester les deux états de l’interrupteur et comparer au schéma. \textbf{À rendre :} nœuds, branches, interrupteur et explication du trajet du courant. Les défauts dangereux sont analysés sur documents, sans les produire.\end{pcEncadre}
\pcSousTitre{Ce que je dois savoir expliquer}\begin{questions}
\item Un schéma utilise des symboles et des connexions explicites.
\item Un courant nécessite un trajet conducteur fermé dans les conditions étudiées.
\item En dérivation, les dipôles sont reliés aux mêmes deux nœuds.
\item Un court-circuit peut provoquer une surintensité.
\item Les expériences utilisent une alimentation scolaire autorisée.
\end{questions}
\begin{pcEncadre}{attention}{Une erreur à éviter}Deux dessins qui se ressemblent géométriquement ne décrivent pas forcément le même circuit : suivre les connexions, pas la position des symboles.\end{pcEncadre}
\pcRepere{cor-activite-13}{Pistes pour l’investigation}\pcGrille
\clearpage\markright{Je démarre}
\rubrique{Je démarre}
\par\noindent\begin{minipage}{\linewidth}
\begin{exercice}[1]{Une boucle}
\pcFigureEx{figures/ex-circuit-ouvert.pdf}{}
\begin{questions}
\item Un interrupteur ouvert dans l’unique boucle permet-il un courant permanent ?
\item Quel est le rôle du générateur ?
\end{questions}
\end{exercice}
\end{minipage}
\par\noindent\begin{minipage}{\linewidth}
\begin{exercice}[1]{Le sens conventionnel}
\begin{questions}
\item Dans le circuit extérieur alimenté en courant continu, de quelle borne part la flèche du courant ?
\item Vers quelle borne revient-elle ?
\end{questions}
\end{exercice}
\end{minipage}
\par\noindent\begin{minipage}{\linewidth}
\begin{exercice}[1]{En série}
\pcFigureEx{figures/ex-circuit-serie.pdf}{}
\begin{questions}
\item Deux lampes sont dans une seule boucle sans dérivation. Comment sont-elles associées ?
\item Que produit l’ouverture d’un interrupteur de cette boucle ?
\end{questions}
\end{exercice}
\end{minipage}
\par\noindent\begin{minipage}{\linewidth}
\begin{exercice}[1]{En dérivation}
\pcFigureEx{figures/ex-circuit-parallele.pdf}{}
\begin{questions}
\item Deux lampes sont reliées aux mêmes deux nœuds. Quelle association ?
\item Peuvent-elles appartenir à deux branches différentes ?
\end{questions}
\end{exercice}
\end{minipage}
\par\noindent\begin{minipage}{\linewidth}
\begin{exercice}[1]{La protection}
\begin{questions}
\item Quel défaut peut provoquer une forte surintensité ?
\item Un fusible intact est-il censé laisser passer indéfiniment toute intensité ?
\end{questions}
\end{exercice}
\end{minipage}
\clearpage\markright{Je m’entraîne}
\rubrique{Je m’entraîne}
\par\noindent\begin{minipage}{\linewidth}
\begin{exercice}[2]{Compter les branches}
\pcFigureEx{figures/ex-circuit-parallele.pdf}{}
\begin{questions}
\item Un générateur et deux lampes sont chacun reliés entre les mêmes deux nœuds. Combien de branches principales au total ?
\item Combien contiennent une lampe ?
\end{questions}
\end{exercice}
\end{minipage}
\par\noindent\begin{minipage}{\linewidth}
\begin{exercice}[2]{Placer l’interrupteur}
\begin{questions}
\item Où placer un interrupteur pour commander seulement la lampe 1 d’un montage en dérivation ?
\item Où le placer pour commander les deux ?
\end{questions}
\end{exercice}
\end{minipage}
\par\noindent\begin{minipage}{\linewidth}
\begin{exercice}[2]{Une lampe dévissée}
\begin{questions}
\item Dans le modèle de deux lampes en dérivation, dévisser L1 ouvre sa branche. Que devient L2 ?
\item Quelle hypothèse faut-il annoncer ?
\end{questions}
\end{exercice}
\end{minipage}
\par\noindent\begin{minipage}{\linewidth}
\begin{exercice}[2]{Suivre les connexions}
\pcFigureEx{figures/ex-circuits-equivalents.pdf}{}
\begin{questions}
\item Deux schémas présentent les mêmes dipôles à des endroits différents. Comment savoir s’ils sont équivalents ?
\item La forme rectangulaire du dessin suffit-elle ?
\end{questions}
\end{exercice}
\end{minipage}
\par\noindent\begin{minipage}{\linewidth}
\begin{exercice}[2]{Un fil parasite}
\pcFigureEx{figures/ex-lampe-court-circuitee.pdf}{Schéma à analyser, sans réaliser le défaut.}
\begin{questions}
\item Un fil relie les bornes d’une lampe. Comment nommer ce défaut ?
\item Que peut-il provoquer dans une boucle série ?
\end{questions}
\end{exercice}
\end{minipage}
\clearpage\markright{Je raisonne}
\rubrique{Je raisonne}
\par\noindent\begin{minipage}{\linewidth}
\begin{exercice}[3]{Un cahier des charges}
\begin{questions}
\item Proposer l’association de deux lampes pour qu’une panne par ouverture de l’une n’interrompe pas l’autre.
\item Cette association garantit-elle contre tous les défauts ?
\end{questions}
\end{exercice}
\end{minipage}
\par\noindent\begin{minipage}{\linewidth}
\begin{exercice}[3]{Un défaut à ne pas tester}
\begin{questions}
\item Pourquoi ne faut-il pas vérifier un court-circuit en reliant une pile par un fil ?
\item Quelle alternative pédagogique convient ?
\end{questions}
\end{exercice}
\end{minipage}
\par\noindent\begin{minipage}{\linewidth}
\begin{exercice}[3]{Une protection et sa limite}
\begin{questions}
\item Un disjoncteur s’est ouvert. Peut-on le réarmer sans rechercher la cause ?
\item Pourquoi une protection ne remplace-t-elle pas l’isolation ?
\end{questions}
\end{exercice}
\end{minipage}
\par\noindent\begin{minipage}{\linewidth}
\begin{exercice}[3]{Une autonomie trompeuse}
\begin{questions}
\item Deux lampes en dérivation brillent, puis s’éteignent ensemble quand on retire un fil commun. Cela contredit-il le fonctionnement indépendant des branches ?
\item Que fallait-il préciser ?
\end{questions}
\end{exercice}
\end{minipage}
\par\noindent\begin{minipage}{\linewidth}
\begin{exercice}[3]{Relier explication et schéma}
\pcFigureEx{figures/ex-circuit-parallele.pdf}{}
\begin{questions}
\item Un élève dessine deux lampes en parallèle mais affirme qu’un seul courant les traverse successivement. Identifier l’incohérence.
\item Comment la corriger ?
\end{questions}
\end{exercice}
\end{minipage}
\pcPage{Défi du chapitre}{Circuits : schématiser et protéger}
\phantomsection\label{defi-ex-13}
\pcDefiTitre{13}{L’éclairage du refuge}
\begin{questions}
\item Le cahier des charges impose deux lampes : L1 fonctionne lorsque l’alimentation est active ; L2 dépend d’un interrupteur. Décrire un schéma.
\item Que se passe-t-il si l’interrupteur de L2 s’ouvre ?
\item Où se situe un défaut qui éteint malgré tout les deux lampes ?
\end{questions}
\pcGrille
